\section*{الفصل \ref{ch:g5:flower-to-fruit} --- من الزهرة إلى الثمرة}

\begin{solution}{exo:g5:flower-to-fruit:1}
\omterm{def:g3:flowers-fruits-seeds:parts}{الأسدية} ذكرية: فحبوب لقاحها تحمل \omterm{def:g5:first-look-at-cells:cell}{الخلايا} التكاثرية الذكرية.
\omterm{def:g3:flowers-fruits-seeds:parts}{والمدقة} أنثوية: ففي داخلها تقع \omterm{prop:g5:flower-to-fruit:roles}{البييضات}، في كل واحدة \omterm{def:g4:animal-reproduction:sexual}{بويضة}.
\end{solution}

\begin{solution}{exo:g5:flower-to-fruit:2}
\omterm{def:g5:flower-to-fruit:steps}{التأبير}: تسافر حبة لقاح من \omterm{def:g3:flowers-fruits-seeds:parts}{سداة} إلى قمة \omterm{def:g3:flowers-fruits-seeds:parts}{مدقة}. \omterm{def:g5:flower-to-fruit:steps}{والإخصاب}: يحمل
أنبوب الحبة \omterm{def:g5:first-look-at-cells:cell}{الخلية} الذكرية نازلًا إلى بييضة، حيث تلتحم \omterm{def:g4:animal-reproduction:sexual}{بالبويضة}.
\omterm{def:g5:flower-to-fruit:steps}{والتأبير} يأتي أولًا --- فالرحلة قبل اللقاء.
\end{solution}

\begin{solution}{exo:g5:flower-to-fruit:3}
تُصنع في \omterm{def:g3:flowers-fruits-seeds:parts}{سداة}؛ وتُحمل (بحشرة أو بريح) إلى قمة \omterm{def:g3:flowers-fruits-seeds:parts}{مدقة} --- وهذا
\omterm{def:g5:flower-to-fruit:steps}{التأبير}؛ ثم تُنبت أنبوبًا دقيقًا نازلًا في \omterm{def:g3:flowers-fruits-seeds:parts}{المدقة}؛ وتبلغ خليتها
الذكرية بييضة فتلتحم \omterm{def:g4:animal-reproduction:sexual}{بالبويضة} --- وهذا \omterm{def:g5:flower-to-fruit:steps}{الإخصاب}.
\end{solution}

\begin{solution}{exo:g5:flower-to-fruit:4}
كل بييضة مخصَّبة تصير \omterm{def:g2:from-seed-to-plant:seed}{بذرة}؛ \omterm{def:g3:flowers-fruits-seeds:parts}{والمدقة} تنتفخ فتصير \omterm{prop:g3:flowers-fruits-seeds:transform}{الثمرة}؛ \omterm{def:g3:flowers-fruits-seeds:parts}{والبتلات}
(\omterm{def:g3:flowers-fruits-seeds:parts}{والأسدية}) تذبل وتسقط.
\end{solution}

\begin{solution}{exo:g5:flower-to-fruit:5}
كل \omterm{def:g2:from-seed-to-plant:seed}{بذرة} بييضة مخصَّبة صارت \omterm{def:g2:from-seed-to-plant:seed}{بذرة}. والحجيرات الفارغة تعني بييضات غير
مخصَّبة --- فقد وقع لقاح أقل من اللازم، فانتفخت \omterm{prop:g3:flowers-fruits-seeds:transform}{الثمرة} انتفاخًا
متفاوتًا حول النقص.
\end{solution}

\begin{solution}{exo:g5:flower-to-fruit:6}
لقاحها يركب الريح: في نعومة الغبار، مصنوعًا بكميات هائلة، مقذوفًا
في الهواء. وثمن ذلك الهدر --- فكل حبة تقريبًا تقع حيث لا فائدة،
فوجب أن تحلّ الكمية محل الدقة.
\end{solution}

\begin{solution}{exo:g5:flower-to-fruit:7}
لا زيارة، فلا لقاح على \omterm{def:g3:flowers-fruits-seeds:parts}{المدقة} --- \omterm{def:g5:flower-to-fruit:steps}{فالتأبير}، وهو الخطوة الأولى، لم
يحدث قط، فلم ينمُ أنبوب ولم تُخصَّب بييضة؛ فتذبل \omterm{def:g3:flowers-fruits-seeds:parts}{الزهرة} كلها.
\end{solution}

\begin{solution}{exo:g5:flower-to-fruit:8}
الجواب مفتوح. والمتوقع: أن تعقد الأزهار المؤبَّرة يدويًّا المعلَّمة
بالخيط ثمرًا أكثر بوضوح --- فالفرشاة أدّت عمل النحلة.
\end{solution}

\begin{solution}{exo:g5:flower-to-fruit:9}
\omterm{def:g5:first-look-at-cells:cell}{الخلايا} الذكرية: في حبوب لقاح \omterm{def:g3:flowers-fruits-seeds:parts}{الأسدية}. \omterm{def:g4:animal-reproduction:sexual}{والبويضات}: في \omterm{prop:g5:flower-to-fruit:roles}{البييضات}
داخل \omterm{def:g3:flowers-fruits-seeds:parts}{المدقة}. \omterm{def:g5:flower-to-fruit:steps}{والإخصاب}: في أسفل أنبوب اللقاح، في البييضة. وأول \omterm{def:g5:first-look-at-cells:cell}{خلية}
للفرد الجديد: \omterm{def:g4:animal-reproduction:sexual}{البويضة} المخصَّبة في البييضة --- وهي التي تنمو \omterm{def:g1:plants-around-us:plant}{نبتة}
\omterm{def:g2:from-seed-to-plant:seed}{البذرة} الصغيرة، كما ينمو \omterm{def:g1:animals-around-us:animal}{الحيوان} من خليته الأولى.
\end{solution}

\begin{solution}{exo:g5:flower-to-fruit:10}
حصاد رديء --- فأكثر الزهر سقط بلا ثمر. وكانت الآلة سليمة والأزهار
سالمة؛ والذي أخفق هو الخطوة الأولى، \omterm{def:g5:flower-to-fruit:steps}{التأبير}: ففي أسبوعين حاسمين
بقي السعاة في بيوتهم، فلم تتلقّ أكثر المدقات لقاحًا ولم تُخصَّب أكثر
\omterm{prop:g5:flower-to-fruit:roles}{البييضات}.
\end{solution}

\begin{solution}{exo:g5:flower-to-fruit:11}
النحل الطنّان يبيع \omterm{def:g5:flower-to-fruit:steps}{التأبير} --- أي خطوة النقل: حمل اللقاح من
\omterm{def:g3:flowers-fruits-seeds:parts}{الأسدية} إلى المدقات داخل بناء لا يستطيع أي مؤبّر خارجي أن يدخله.
ولقاح الطماطم ينفرط بسهولة، فتستطيع اهتزازة مروحة وحركة الهواء أن
تقوما بنقل خشن داخل كل \omterm{def:g3:flowers-fruits-seeds:parts}{زهرة} --- ساعيًا آليًّا. أما لقاح الكرز فعليه
أن يسافر من شجرة إلى شجرة ومن \omterm{def:g3:flowers-fruits-seeds:parts}{زهرة} إلى \omterm{def:g3:flowers-fruits-seeds:parts}{زهرة}، بدقة لا يوفّرها إلا
جسم يزور الأزهار واحدة بعد أخرى؛ والريح في قاعة لا تصوّب، فليس
للساعي الحي في الكرز بديل رخيص.
\end{solution}
